% ============================================================================
%  BUAS BACHELOR'S THESIS TEMPLATE
%  Academy of Games and Media — Data Science & AI
%  
%
%  INSTRUCTIONS:
%  1. Replace all placeholder text (marked with <<<...>>>) with your content
%  2. Read the guidance in the teal boxes for section-specific advice
%  3. Read the grey comment blocks for formatting and style tips
%  4. Delete ALL \guidancebox and \stylenote environments before submission
%  5. Compile with: pdflatex thesis.tex && bibtex thesis && pdflatex thesis.tex && pdflatex thesis.tex
%
%  Academic Year 2025-2026
% ============================================================================

\documentclass[12pt,a4paper,twocolumn]{article}

% ── Page geometry ────────────────────────────────────────────────────────
% A4 paper, with 25mm margins on all sides. The title, subtitle, author
% information, and abstract are presented in a single column, while the
% main text is formatted in two columns. Each column is 77mm wide, with
% 6mm of space between them.
\usepackage[
  a4paper,
  left=25mm,
  right=25mm,
  top=25mm,
  bottom=25mm,
  columnsep=6mm
]{geometry}

% ── Fonts ────────────────────────────────────────────────────────────────
% The font should be Palatino Linotype (Regular), or a similar font
% (e.g. Computer Modern in LaTeX). Times Roman and Georgia are NOT
% recommended, as they are less suitable for long academic texts.
% Use 10pt font with single line spacing.
% Use bold only for subheadings. Use italics for emphasis, foreign words
% (except commonplace terms), Latin abbreviations, and newly introduced
% terminology. NEVER use underlining — it belongs to the era of typewriters.
\usepackage[T1]{fontenc}
\usepackage[utf8]{inputenc}
\usepackage{mathpazo}             % Palatino text + math
\usepackage{microtype}            % Improved typographic quality
\linespread{1.0}                  % Single line spacing

% ── Essential packages ───────────────────────────────────────────────────
\usepackage{amsmath,amssymb}      % Mathematics
\usepackage{graphicx}             % Images
\usepackage{booktabs}             % Professional tables (\toprule, \midrule, \bottomrule)
\usepackage{caption}              % Caption formatting
\usepackage{float}                % Improved float placement
\usepackage{enumitem}             % List customization
\usepackage{url}                  % URL formatting
\usepackage{algorithm2e}          % Pseudo code
\usepackage{xcolor}               % Colours
\usepackage{fancyhdr}             % Headers and footers
\usepackage{titlesec}             % Section heading formatting
\usepackage{natbib}               % APA-style citations
\usepackage{hyperref}             % Clickable links (load last)

% ── Colour definitions (BUas palette) ────────────────────────────────────
\definecolor{buasnavy}{HTML}{0F2A44}
\definecolor{buasteal}{HTML}{0D7377}
\definecolor{buasgold}{HTML}{E8A838}
\definecolor{buascream}{HTML}{F7F3ED}
\definecolor{guidegray}{HTML}{5D6D7E}
\definecolor{guideboxbg}{HTML}{E8F6F6}
\definecolor{styleboxbg}{HTML}{F5F5F0}
\definecolor{warnred}{HTML}{C0392B}

% ── Hyperlink styling ────────────────────────────────────────────────────
\hypersetup{
  colorlinks=true,
  linkcolor=buasnavy,
  citecolor=buasteal,
  urlcolor=buasteal,
  pdfauthor={<<<Your Name>>>},
  pdftitle={<<<Title of Your Thesis>>>},
}

% ── Caption formatting ───────────────────────────────────────────────────
% Figures: caption BELOW in 9pt bold — "Figure X.Y: Description"
% Tables: caption ABOVE in 9pt bold — "Table X.Y: Description"
% Leave 6pt between caption and content.
\captionsetup{
  font={bf,small},
  labelsep=colon,
  skip=6pt,
}
\captionsetup[table]{position=above}
\captionsetup[figure]{position=below}

% ── Section heading formatting ───────────────────────────────────────────
% Level 1: 12pt bold, 12pt space before, 3pt after. Arabic numerals (1, 2, 3).
% Level 2: 10pt bold, 12pt space before, 3pt after. Double numbers (1.1, 1.2).
% Subheadings begin with a capital letter but contain no additional capitals.
% They do not end with a full stop. Use a maximum of two levels in a short text.
\titleformat{\section}
  {\normalfont\fontsize{12}{14}\bfseries}{\thesection}{1em}{}
\titlespacing*{\section}{0pt}{12pt}{3pt}

\titleformat{\subsection}
  {\normalfont\fontsize{10}{12}\bfseries}{\thesubsection}{1em}{}
\titlespacing*{\subsection}{0pt}{12pt}{3pt}

% ── Page numbering ───────────────────────────────────────────────────────
% Page numbers appear at the bottom right of the page, 10mm below the text.
\pagestyle{fancy}
\fancyhf{}
\fancyfoot[R]{\small\thepage}
\renewcommand{\headrulewidth}{0pt}
\renewcommand{\footrulewidth}{0pt}

% ── Paragraph formatting ─────────────────────────────────────────────────
% Paragraphs should be justified. Words may be hyphenated.
% Do not indent the first paragraph of a section. Subsequent paragraphs
% should have a 6mm first-line indent.
\setlength{\parindent}{6mm}
\setlength{\parskip}{0pt}

% ── Algorithm styling ────────────────────────────────────────────────────
% Pseudo code at 9pt. Variables in italics, control words (if, then, while)
% in boldface, function names in small caps. Use line numbers if you
% reference specific lines. Blocks indented, closed with end for, etc.
% Caption ABOVE the code: "Algorithm X.Y: description" in 9pt bold.
\SetAlgorithmName{Algorithm}{Algorithm}{List of Algorithms}
\SetAlgoLined
\SetKwComment{Comment}{/* }{ */}

% ── Guidance box environments ────────────────────────────────────────────
% TEAL BOXES: Section-specific writing guidance (what to write).
% GREY BOXES: Style and formatting notes (how to write and format).
% >>> DELETE ALL OF THESE BEFORE FINAL SUBMISSION <<<
\usepackage{tcolorbox}
\tcbuselibrary{breakable}

% Content guidance (teal)
\newtcolorbox{guidancebox}[1][]{
  colback=guideboxbg,
  colframe=buasteal,
  fonttitle=\bfseries\small\color{buasteal},
  title={Writing Guidance — Delete before submission},
  coltitle=buasteal,
  boxrule=0.5pt,
  arc=2pt,
  left=4pt, right=4pt, top=4pt, bottom=4pt,
  fontupper=\small\color{guidegray},
  breakable,
  before skip=6pt, after skip=6pt,
  #1
}

% Style/formatting notes (grey)
\newtcolorbox{stylenote}[1][]{
  colback=styleboxbg,
  colframe=guidegray,
  fonttitle=\bfseries\small\color{guidegray},
  title={Style \& Formatting Note — Delete before submission},
  coltitle=guidegray,
  boxrule=0.5pt,
  arc=2pt,
  left=4pt, right=4pt, top=4pt, bottom=4pt,
  fontupper=\small\color{guidegray},
  breakable,
  before skip=6pt, after skip=6pt,
  #1
}

% ════════════════════════════════════════════════════════════════════════════
%  DOCUMENT BEGINS
% ════════════════════════════════════════════════════════════════════════════
\begin{document}

% ── Title block (single column, full width) ──────────────────────────────
% The title, subtitle, author information, and abstract span the full page width.
\twocolumn[
  \begin{@twocolumnfalse}

  % ── University / programme header ──────────────────────────────────────
  \begin{center}
    {\small\color{guidegray}
      Breda University of Applied Sciences \hspace{1em}$\mid$\hspace{1em}
      Academy of Games and Media \hspace{1em}$\mid$\hspace{1em}
      Data Science \& AI}
    \vspace{12pt}

    % ── Title ────────────────────────────────────────────────────────────
    % 16pt bold, small caps for lowercase letters. Content words begin
    % with a capital letter ('title caps'). Leave 3pt below the title.
    {\fontsize{16}{19}\bfseries\scshape
      <<<Title of Your Thesis Project>>>
    \par}
    \vspace{3pt}

    {\normalsize Bachelor's Thesis}
    \vspace{6pt}

    % ── Author information ───────────────────────────────────────────────
    % Name, student number, email. Remove the student number from the
    % final version you upload to the repository.
    {\normalsize
      <<<Your Name>>>, <<<Student Number>>>, <<<your.email@buas.nl>>>}
    \vspace{2pt}

    {\normalsize
      Supervisors: <<<First Supervisor>>> \& <<<Second Supervisor>>>}
    \vspace{12pt}

    % ── Abstract ─────────────────────────────────────────────────────────
    % 9pt font, single column across the full page, with a 6mm indent on both
    % sides. Maximum 250 words. Begins with "Abstract:" in bold.
    % The abstract is separate from the rest of the text — the main text
    % should be understandable without first reading the abstract.
    \begin{minipage}{0.92\textwidth}
      \small
      \setlength{\parindent}{6mm}

      \textbf{Abstract:}
      <<<Write your abstract here. Maximum 250 words. The abstract serves as
      the ``teaser'' of your thesis. On the basis of the abstract, readers
      decide whether the rest of the thesis is worth reading. In many
      databases, the abstract is the only freely available part.

      The abstract MUST include: (1)~the research problem and its context,
      (2)~your approach or method, (3)~the key findings and quantitative
      results, and (4)~the main conclusions.  It should provide a complete picture of what was done and what can be concluded from it.

      The abstract must be readable as a stand-alone text: do not refer to
      sections, figures, or tables. Write the abstract LAST, once your
      results and conclusions are clear.>>>

    \end{minipage}

    \vspace{6pt}
  \end{center}

  \end{@twocolumnfalse}
]

% ════════════════════════════════════════════════════════════════════════════
%  SECTION 1: INTRODUCTION
% ════════════════════════════════════════════════════════════════════════════
\section{Introduction}
\label{sec:introduction}

\begin{guidancebox}
\textbf{Purpose:} The introduction and conclusion together should provide a
\emph{complete and comprehensible} picture of your thesis. This is intended
for time-poor readers whose interest has been raised by the abstract but
who do not have time to read the full thesis.

\textbf{Your introduction must:}
\begin{itemize}[leftmargin=*, nosep]
  \item Set out the \textbf{context} of the work: why it is relevant and
    who has previously worked on the topic.
  \item State the \textbf{specific research question}, and explain how it
    contributes to a better understanding of the subject.
  \item Provide a \textbf{summary of the techniques and analysis methods}
    used (details are given in the Methodology section).
  \item \textbf{Convince the reader} that your work is useful and
    interesting, that you understand the subject well, and that you have
    chosen an appropriate way to investigate it.
\end{itemize}



\textbf{Delete this box before submission.}
\end{guidancebox}

\begin{stylenote}
\textbf{Before you start writing}, answer these four questions:
\begin{enumerate}[leftmargin=*, nosep]
  \item \textbf{What is the aim?} Are you reporting on a project and its
    results? Presenting a new idea? Proposing a solution? For a BSc thesis,
    the aim is usually to report on a research project and present your findings.
  \item \textbf{Who is your audience?} Your supervisor, industry readers,
    and fellow researchers are the primary audience. They are knowledgeable
    but critical. Do not over-explain basic concepts, but substantiate every claim.
  \item \textbf{What is the format?} A BSc thesis is typically 10--15 pages
    (excluding references and appendices). This is closer to a conference paper than to a
    journal article, so scope your depth and level of detail accordingly.
  \item \textbf{What is your contribution?} State your specific
    contribution clearly: what message are you trying to convey? What can
    you assume is already known, and what background do you need to provide?
\end{enumerate}

\textbf{Tip:} Prepare a detailed table of contents before you begin writing
chapters and sections, and add a brief description of each. This helps
prevent structural problems later.

\textbf{Delete this box before submission.}
\end{stylenote}

\noindent
<<<Write your introduction here. Start by establishing the broader context
of your work, then narrow the focus to your specific research question.
Provide a brief overview of your approach and convince the reader that your
work is both necessary and well motivated \cite{brown2020}.>>>


% ════════════════════════════════════════════════════════════════════════════
%  SECTION 2: RELATED WORK
% ════════════════════════════════════════════════════════════════════════════
\section{Related Work}
\label{sec:related-work}

\begin{guidancebox}
\textbf{Purpose:} Position your research within the existing body of
knowledge. This section is \emph{not} a literature dump. It should present
a structured argument for why your research is necessary. Every paper you
cite should help build that argument.

\textbf{What to include:}
\begin{itemize}[leftmargin=*, nosep]
  \item Prior approaches to your problem or similar problems.
  \item Key results and findings from the literature.
  \item Theoretical foundations your work builds upon.
  \item Gaps in the literature that your work addresses.
\end{itemize}

\textbf{How to structure:}
\begin{itemize}[leftmargin=*, nosep]
  \item Group the literature by \textbf{theme or method}, not chronologically.
  \item \textbf{Relate each cited work} to your own research question.
    Do not simply summarise it; explain why it matters.
  \item Move from the \textbf{broader context to the specific gap} that your work fills.
  \item End by clearly stating how your work \textbf{extends or differs}
    from prior research.
\end{itemize}

\textbf{Pitfalls:}
\begin{itemize}[leftmargin=*, nosep]
  \item Listing papers without connecting them to your work.
  \item Summarising the literature without critical analysis or comparison.
  \item Writing literature overviews that summarise a large amount of prior
    work without clearly relating it to the author's own study.
  \item Relying only on textbooks or secondary sources.
  \item Overlooking recent publications (last 1--2 years).
\end{itemize}

\textbf{Test:} After reading this section, the reader should be able to explain
(a)~what has already been tried, (b)~what worked and what did not, and
(c)~exactly which gap your thesis fills.

\textbf{Delete this box before submission.}
\end{guidancebox}

\noindent
<<<Write your related work here. Group the literature thematically, relate
each reference to your research question, and clearly articulate the gap
your work addresses \cite{dawkins1976,brown2020}.>>>


% ════════════════════════════════════════════════════════════════════════════
%  SECTION 3: METHODOLOGY
% ════════════════════════════════════════════════════════════════════════════
\section{Methodology}
\label{sec:methods}

\begin{guidancebox}
\textbf{Golden rule:} The description must focus on \textbf{reproducibility}.
A reader should be able to replicate your work on the basis of this section alone.
For a computer program, sufficient detail must be provided to allow the
implementation to be repeated. For an experiment, sufficient detail must
be given to allow the procedure to be reproduced accurately.

\textbf{Three pillars:}
\begin{enumerate}[leftmargin=*, nosep]
  \item \textbf{Describe what, not how}: explain the system above the
    source-code level. Describe what each component does, rather than how
    it was programmed. A good approach is to first describe how the system
    works as a whole, and then discuss the individual components.
    Substantiate the description with formulae and pseudo code.
  \item \textbf{Justify your choices}: in many cases, aspects of your
    system can be implemented in different ways. Explain why you selected
    a particular method \emph{and why you rejected the alternatives}.
    Place these choices in the broader context of your research question.
    Describing the system involves more than simply stating what you built.
  \item \textbf{Include your analysis plan}: the method section should
    also describe the techniques used to \emph{analyse} the results.
    Define the evaluation metrics, describe how raw data were collected,
    averaged, and processed, and name the statistical tests used while
    justifying their validity.
\end{enumerate}

\textbf{Important:} Present all parameters in a central table (see
Table~\ref{tab:parameters}). The reader should never have to search for a
value buried in the prose.

\textbf{Delete this box before submission.}
\end{guidancebox}

\noindent
<<<Describe your methodology here. Start with a general overview, then
describe the individual components. Justify every design decision and
include your full analysis plan \cite{devlin2019}.>>>

% ── Example: parameter table ─────────────────────────────────────────────
\begin{table}[t]
  \caption{System parameters.}
  \label{tab:parameters}
  \centering
  \small
  \begin{tabular}{@{}ll@{}}
    \toprule
    \textbf{Parameter}      & \textbf{Value} \\
    \midrule
    Learning rate            & <<<0.001>>>    \\
    Batch size               & <<<32>>>       \\
    Hidden dimensions        & <<<256>>>      \\
    Number of epochs         & <<<50>>>       \\
    Dropout rate             & <<<0.3>>>      \\
    Optimizer                & <<<Adam>>>     \\
    <<<\ldots>>>             & <<<\ldots>>>   \\
    \bottomrule
  \end{tabular}
\end{table}

% ── Example: algorithm ───────────────────────────────────────────────────
% Pseudo code: 9pt, set apart from the main text. Variables in italics,
% control words (if, then, while) in boldface, function names (Sort,
% FindShortest) in small caps. Line numbers required if referencing
% specific lines. Blocks are indented and closed with end for / end while.
% Caption ABOVE: "Algorithm X.Y: description" in 9pt bold.
% In-text: "see Algorithm~\ref{alg:example}".
\begin{algorithm}[t]
  \caption{<<<Name of your algorithm>>>}
  \label{alg:example}
  \KwIn{<<<Input description>>>}
  \KwOut{<<<Output description>>>}
  <<<Replace with your algorithm in pseudo code>>>\;
  \For{<<<each item in dataset>>>}{
    <<<Process item>>>\;
    \If{<<<condition>>>}{
      <<<Update model>>>\;
    }
  }
  \Return{<<<result>>>}\;
\end{algorithm}


% ════════════════════════════════════════════════════════════════════════════
%  SECTION 4: RESULTS
% ════════════════════════════════════════════════════════════════════════════
\section{Results}
\label{sec:results}

\begin{guidancebox}
\textbf{Purpose:} Present the essence of your message: what you
discovered and how your findings shed light on the research question.
Two things are especially important: \textbf{selection} and \textbf{presentation}.

\textbf{Selection:} Present results that answer your research question.
If a result does not strengthen your argument, you may leave it out.

\textbf{Presentation format:}
\begin{itemize}[leftmargin=*, nosep]
  \item \textbf{Line graphs} for time series or continuous relationships.
  \item \textbf{Bar charts} for categorical comparisons.
  \item \textbf{Tables} when precise numerical values matter.
  \item \textbf{Prose} when qualitative interpretation is the main point.
\end{itemize}

\textbf{Statistical significance:}
\begin{itemize}[leftmargin=*, nosep]
  \item Graphs: use error bars, box plots, or confidence intervals to
    show variability.
  \item Tables: include standard deviations, confidence intervals, and
    p-values.
  \item Text: name the statistical tests (t-test, ANOVA, Wilcoxon, etc.)
    and justify their validity.
\end{itemize}

\textbf{Caution:} Do not confuse significance with relevance. A
statistically significant result can still be practically irrelevant.

\textbf{Delete this box before submission.}
\end{guidancebox}

\noindent
<<<Present your results here. Select only findings that are relevant to
your research question and present them in the most appropriate format.
Always include measures of statistical significance.>>>

% ── Example: results table ───────────────────────────────────────────────
% Tables: 9pt font, with rows/columns separated by single half-point lines.
% Labels in bold. Numbers centred or aligned on the decimal point.
% Caption ABOVE: "Table X.Y: description" in 9pt bold.
% In-text: "see Table~\ref{tab:results}".
\begin{table}[t]
  \caption{<<<Description of your results table.>>>}
  \label{tab:results}
  \centering
  \small
  \begin{tabular}{@{}lccc@{}}
    \toprule
    \textbf{Model}  & \textbf{Precision} & \textbf{Recall} & \textbf{F1} \\
    \midrule
    Baseline         & <<<0.72>>>         & <<<0.68>>>      & <<<0.70>>>  \\
    Your model       & <<<0.85>>>         & <<<0.81>>>      & <<<0.83>>>  \\
    \bottomrule
  \end{tabular}
\end{table}

% ── Example: figure ──────────────────────────────────────────────────────
% Figures: use vector graphics (PDF, EPS) for diagrams. Bitmaps (PNG, JPG)
% ONLY for photographs, and only if the resolution is high enough that
% there is no visible quality loss.
% Graphs: use clear axes with marked values, indicate what they show,
% and include units where relevant. If multiple items are shown, include a legend.
% All elements must remain clearly distinguishable even in BLACK AND WHITE.
% Include reliability intervals or other indicators of variation.
% Figures may be one or two columns wide and should be placed within the margins.
% Caption BELOW: "Figure X.Y: description" in 9pt bold. Leave 6pt below the caption.
% In-text: "see Figure~\ref{fig:example}".
%
% \begin{figure}[t]
%   \centering
%   \includegraphics[width=\columnwidth]{figures/<<<your_figure>>>}
%   \caption{<<<Description. Ensure that axes and labels are clearly legible.>>>}
%   \label{fig:example}
% \end{figure}


% ════════════════════════════════════════════════════════════════════════════
%  SECTION 5: DISCUSSION
% ════════════════════════════════════════════════════════════════════════════
\section{Discussion}
\label{sec:discussion}

\begin{guidancebox}
\textbf{Purpose:} Describe the implications of your findings in the
context of your research question. Guide the reader through the results by
explaining how each finding relates to the research question and to prior
work in the field. This discussion should build the argument for a short
and, ideally, powerful final conclusion.

\textbf{What to cover:}
\begin{itemize}[leftmargin=*, nosep]
  \item \textbf{Interpretation:} explain how each result relates to your
    research question and guide the reader through that interpretation.
  \item \textbf{Comparison with related work:} do your findings align with,
    contradict, or extend prior results? Link back to
    Section~\ref{sec:related-work}.
  \item \textbf{Broader context:} place your conclusions in the wider
    academic context of the field.
  \item \textbf{Practical implications:} what do your findings mean for
    practitioners?
  \item \textbf{Future research:} suggest concrete future research
    directions that follow from your findings.
\end{itemize}

\textbf{Limitations \& pitfalls:}
\begin{itemize}[leftmargin=*, nosep]
  \item Acknowledge limitations honestly; reviewers will identify them anyway.
  \item Discuss threats to validity (internal, external, construct).
  \item Do not assume that adding complexity is the default solution;
    simplification may sometimes be more effective.
  \item Do not confuse statistical significance with practical relevance.
  \item Future work should follow from your analysis, not read like a wish list.
\end{itemize}

\textbf{Delete this box before submission.}
\end{guidancebox}

\noindent
<<<Interpret your results here. Explain what they mean for the research
question, compare them with related work, acknowledge limitations honestly,
and suggest concrete directions for future research \cite{vaswani2017}.>>>


% ════════════════════════════════════════════════════════════════════════════
%  SECTION 6: CONCLUSION
% ════════════════════════════════════════════════════════════════════════════
\section{Conclusion}
\label{sec:conclusion}

\begin{guidancebox}
\textbf{Purpose:} The introduction and conclusion together should form a
\emph{mini-article} that can be read on its own. Include all of your key
findings in this section. There is no need to repeat graphs and tables,
but it is essential to summarise your findings in words.

\textbf{Must include:}
\begin{itemize}[leftmargin=*, nosep]
  \item \textbf{Summarize key findings in words}: the reader should
    understand your results without needing to read the full thesis.
  \item \textbf{Directly answer the research question} introduced earlier.
    This is the single most important paragraph.
  \item \textbf{Briefly restate your approach} so that the conclusion can
    stand on its own.
  \item \textbf{Connect back to related work}: does your study confirm,
    extend, or contradict prior research?
  \item \textbf{Be concise}: state each conclusion clearly and separately.
\end{itemize}

\textbf{Self-test:} Give someone only your Introduction and Conclusion.
If they can explain what you did, why you did it, what you found, and what
it means, then your thesis structure is working well.

\textbf{Delete this box before submission.}
\end{guidancebox}

\noindent
<<<Write your conclusion here. Summarize your findings, directly answer your
research question, and connect your conclusions back to the broader context
established in the introduction.>>>


% ════════════════════════════════════════════════════════════════════════════
%  ACADEMIC WRITING STYLE GUIDANCE
% ════════════════════════════════════════════════════════════════════════════
% This box appears between the conclusion and references. It contains
% guidance only — delete it entirely before submission.

\begin{stylenote}[title={Academic Writing Style Guide — Delete before submission}]

\textbf{DO:}
\begin{itemize}[leftmargin=*, nosep]
  \item Use impersonal constructions or passive voice where appropriate.
  \item ``This study aimed to test\ldots'' rather than ``I wanted to see\ldots''
  \item Substantiate every claim with references or data.
  \item Build paragraphs: topic sentence $\rightarrow$ detail $\rightarrow$ transition.
  \item Ask a peer to read and critique your draft.
  \item Read published papers in your field to learn the conventions.
\end{itemize}

\textbf{DON'T:}
\begin{itemize}[leftmargin=*, nosep]
  \item Overuse first-person pronouns. Avoid ``we'' in a single-author thesis.
  \item Make vague generalizations such as ``Many researchers believe\ldots''.
    Cite them or remove them.
  \item Use overly formal language that reduces readability.
  \item Include material that does not serve your main message.
  \item Use Wikipedia as a primary source. Use it only to locate proper sources.
\end{itemize}

\medskip
\textbf{Detailed guidance below:}

\textbf{The academic style.}
Academic writing generally avoids first- and second-person pronouns
(I, you, we). Use impersonal constructions or passive voice where appropriate.
Write ``this study aimed to test\ldots'' rather than ``I wanted to see\ldots''.
In a single-author thesis, ``we'' is usually inappropriate rather than helpful.

\textbf{Paragraph construction.}
A paragraph should focus on one main subject, or link two closely related ones.
The first sentence introduces the topic; the final sentence provides a transition
to the next paragraph; the middle sentences provide the supporting detail.
If you cannot answer ``Is this relevant to my message?'' positively, leave it out.
If you cannot answer ``Will my readers understand this?'' positively, explain it in more detail.

\textbf{Substantiate claims.}
Vague generalizations such as ``The interest in X is rapidly increasing'' or
``Many scientists believe that\ldots'' should be avoided unless they are
properly substantiated with references.

\textbf{Get feedback.}
Writers are often the least critical readers of their own work. Ask someone
else to read your text and take their comments seriously. If \emph{you}
already think something is weak, it is very likely to be weak.

\textbf{Numbers and notation.}
Use modified American notation: points for decimals (1.5, not 1,5).
Scientific notation for large numbers: $6.022 \times 10^{23}$.
Use spaces to group digits: 4\,294\,967\,296. Do not allow line breaks in
the middle of formulae or between a number and its unit (e.g., 6\,cm).

\textbf{Formulae.}
Use standard mathematical notation. Explain all symbols clearly. Formulae
should be left-aligned and their numbers right-aligned (e.g., Equation~4.1).
Use the same font as in the main text.

\textbf{Italics.}
Use italics for words from other languages (except commonplace terms),
Latin abbreviations, new terminology, emphasis, and short quotations.
Never underline.

\textbf{Dashes.}
Hyphen (-): joins words or separates syllables.
En-dash (--): minus sign, number ranges ($-1$, 2--5).
Em-dash (---): pause or parenthetical.

\textbf{Quotation marks.}
Use proper opening and closing marks: `correct' and ``also correct''.
Not 'wrong' or "wrong again".

\textbf{Spelling.}
Use a spelling checker. Check that accents and special characters are
correct. Verify hyphenation, as LaTeX may occasionally make mistakes.

\textbf{Read widely.}
One of the most useful ways to improve your writing is to read many
academic articles and observe what works and what does not. A good article
will always remain a good article, and a bad one will remain a bad one.
Reading both helps you understand why certain writing choices are effective.

\textbf{Delete this entire box before submission.}
\end{stylenote}


% ════════════════════════════════════════════════════════════════════════════
%  REFERENCES
% ════════════════════════════════════════════════════════════════════════════
% APA style. In-text citation patterns:
%   One author:        \citep{dawkins1976}            → (Dawkins, 1976)
%   Two authors:       \citep{berrah1999}             → (Berrah & Laboissière, 1999)
%   Three or more:     \citep{schwartz1997}           → (Schwartz et al., 1997)
%   Multiple sources:  \citep{cooper1952,crothers1975} → (Cooper et al., 1952; Crothers, 1975)
%   Same author, multi-year: \citep{kirby1998,kirby1999,kirby2000}
%   Secondary source:  (Schwartz et al., 1997, as cited in Berrah & Laboissière, 1999)
%     — list both in the bibliography.
%   Narrative:         \citet{dawkins1976}            → Dawkins (1976)
%
% WHEN TO REFERENCE:
%   - Whenever you use other people's ideas, text, or illustrations.
%   - When describing system choices (must be based on what others tried).
%   - In the introduction: similar work, without interrupting the flow.
%   - In the conclusion: similar conclusions by fellow researchers.
%
% WHICH SOURCE?
%   - Refer to both the original discoverer AND the clearer textbook.
%   - Better to over-reference than under-reference.
%   - If you haven't read the original, be open:
%     (Schwartz et al., 1997, as cited in Berrah & Laboissière, 1999)
%   - Avoid vague, hard-to-find, in-press, or online-only sources.
%
% The bibliography must include: author, year, title, publisher/journal,
% and issue number. Use APA conventions (or another style approved by
% your supervisor beforehand).

\newpage
\onecolumn
\bibliographystyle{apalike}
\bibliography{references}

% ── Fallback: manual bibliography (uncomment if no .bib file) ───────────
% \begin{thebibliography}{99}
%
%   \bibitem[Dawkins(1976)]{dawkins1976}
%     Dawkins, R. (1976). \textit{The Selfish Gene}. Oxford University Press.
%
%   \bibitem[Vaswani et~al.(2017)]{vaswani2017}
%     Vaswani, A., Shazeer, N., Parmar, N., et~al. (2017).
%     Attention is all you need.
%     \textit{Advances in Neural Information Processing Systems}, 30.
%
%   \bibitem[Devlin et~al.(2019)]{devlin2019}
%     Devlin, J., Chang, M.-W., Lee, K., \& Toutanova, K. (2019).
%     {BERT}: Pre-training of deep bidirectional transformers for language
%     understanding. \textit{Proceedings of NAACL-HLT}, 4171--4186.
%
% \end{thebibliography}


% ════════════════════════════════════════════════════════════════════════════
%  APPENDICES
% ════════════════════════════════════════════════════════════════════════════
% Appendices are NOT counted towards the 10–15 page limit.
% Use letters (A, B, C) to number them. You may use single-column layout.
% Start the first appendix on a new page.
% If you have many appendices, provide a list on the first page.

\newpage
\onecolumn  % Optional: single-column layout for appendices

\appendix

\section{<<<Title of First Appendix>>>}
\label{app:first}

<<<Include supplementary material here: extended tables, additional
figures, code listings, detailed derivations, or other material that
did not fit in the main text but supports reproducibility. If you
have many appendices, provide a list at the beginning.>>>

% Uncomment for additional appendices:
% \section{<<<Title of Second Appendix>>>}
% \label{app:second}
% <<<Additional supplementary material.>>>


\end{document}
